%It has been shown that constraining the Lipschitz constant of a deep neural networks increases the robustness against adversarial attacks. We propose a new framework of binary classification based on optimal transport that integrates these structural constraint as a theoretical requirement. We propose to compute networks that are the solution of an hinge regularized version of the Kantorovich-Rubinstein dual formulation for the estimation of the Wasserstein distance. The loss function obtained has a direct interpretation in terms of adversarial robustness together with certifiable robustness bound. We prove that this classification formulation has a solution, and is still the dual formulation of an optimal transportation problem despite the use of the hinge regularization. We also establish the geometrical properties of this optimal solution and we extend the approach to multi-class problems. The experiments show that the approach provides the expected guarantees in terms of robustness without any significant accuracy drop. The results also suggest that adversarial attacks on the proposed models visibly and meaningfully change the input, and can thus serve as an explanation for the classification.

Adversarial examples have pointed out Deep Neural Networks vulnerability to small local noise. It has been shown that constraining their Lipschitz constant should enhance robustness, but make them harder to learn with classical loss functions. We propose a new framework for binary classification, based on optimal transport, which integrates this Lipschitz constraint as a theoretical requirement. We propose to learn 1-Lipschitz networks using a new loss that is an hinge regularized version of the Kantorovich-Rubinstein dual formulation for the Wasserstein distance estimation. This loss function has a direct interpretation in terms of adversarial robustness together with certifiable robustness bound. We also prove that this hinge regularized version is still the dual formulation of an optimal transportation problem, and has a solution. We also establish several geometrical properties of this optimal solution, and extend the approach to multi-class problems. Experiments show that the proposed approach provides the expected guarantees in terms of robustness without any significant accuracy drop. The adversarial examples, on the proposed models, visibly and meaningfully change the input providing an explanation for the classification.



% We also experimented adversarial attack on the proposed model, which suggests that visible and meaningfull changes must be made to the input in order to be mis-classified. This opens perspectives to 
%The results also suggest that adversarial attacks on the proposed models require to perform explicit change on the input critical part, and can thus serve as an explanation for the classification.


